\documentclass[12pt,a4paper,oneside]{article}
\usepackage[brazil]{babel}
\usepackage[utf8]{inputenc}
\usepackage{olymp}

\setcounter{problem}{2}

\begin{document}

\begin{problem}{Siglas [2]}{siglas2}{2}
 
Um acrônimo, ou mais comumente conhecido como sigla, é uma abreviação formada pelos componentes iniciais de uma palavra ou de uma frase. Estes componentes podem ser letras individuais, como em UFV (\underline{\textbf{U}}niversidade \underline{\textbf{F}}ederal
de \underline{\textbf{V}}içosa), ou partes iniciais das palavras, como em ANVISA (\underline{\textbf{A}}gência \underline{\textbf{N}}acional de \underline{\textbf{Vi}}gilância \underline{\textbf{Sa}}nitária).

Na verdade, não existe um padrão muito rígido para a formação de siglas, e existem siglas que contém partes das palavras que não estão necessariamente no início delas, como em DPI (\underline{\textbf{D}}e\underline{\textbf{p}}artamento de \underline{\textbf{I}}nformática) ou em CCP (\underline{\textbf{C}}iência da \underline{\textbf{C}}om\underline{\textbf{p}}utação).

Algumas siglas parecem até não fazer sentido, como SAPIENS (Sistema de Apoio
ao Ensino), mas após uma análise é possível perceber que todas as letras
utilizadas na sigla estão presentes na frase original, e na ordem utilizada (\underline{\textbf{S}}istema
de \underline{\textbf{Ap}}o\underline{\textbf{i}}o ao \underline{\textbf{Ens}}ino). Mas existem ainda outras que realmente não fazem sentido, já que não é possível encontrar todas as letras da sigla na frase original, na
mesma ordem. Vários exemplos podem ser encontrados na lista de códigos dos
aeroportos, como BSB (Brasília - DF) e VIX (Vitória – ES).

Faça um programa que, dada uma ou mais palavras e uma sigla correspondente,
determine se aquela sigla faz sentido ou não. Para que uma sigla faça sentido, é
necessário que todas as letras dela estejam presentes na frase original, na
mesma ordem.

%\Notes
%
%Observações, se tiver.

\InputFile

A entrada contém duas linhas. A primeira linha contém uma única palavra, indicando a sigla utilizada, e a segunda contém uma ou mais palavras, indicando a frase original. Todas as letras da entrada são letras maiúsculas, e as palavras da frase são separadas por espaços. Cada linha da entrada terá no máximo 200
caracteres.

\OutputFile

Seu programa deve escrever somente ``\texttt{SIM}'' ou ``\texttt{NAO}'', indicando se a sigla dada faz sentido ou não para aquela frase.

% \Notes
% Preste atenção aos tipos das variáveis, pois $F(90) \approx 2^{61}$.

\Example

\begin{example}
\exmp{
DPI
DEPARTAMENTO DE INFORMATICA
}{
SIM
}%
\end{example}

\begin{example}
\exmp{
SAPIENS
SISTEMA DE APOIO AO ENSINO
}{
SIM
}%
\end{example}

\begin{example}
\exmp{
BSB
BRASILIA DF
}{
NAO
}%
\end{example}

\begin{example}
\exmp{
VIX
VITORIA ES
}{
NAO
}%
\end{example}

\begin{example}
\exmp{
BBB
BENSON MUNICIPAL AIRPORT MN
}{
NAO
}%
\end{example}

%Comentários sobre o exemplo, se necessário.

\end{problem}

\end{document}
